\pagebreak

\section{AVNET Microzed Evaluation Board Header File - boards/avnet/microzed\_70x0.3pl}

    
    These library must be incorporated by using a {\bf module} statements -
    \begin{lstlisting}[gobble=4, language=threepl]
       module "boards/avnet/microzed_7010.3pl"
    \end{lstlisting}
    or
    \begin{lstlisting}[gobble=4, language=threepl]
       module "boards/avnet/microzed_7020.3pl"
    \end{lstlisting}
    
    The following is included -
    \begin{lstlisting}[gobble=4, language=threepl]
       source "boards/avnet/microzed_base.3pl"
    \end{lstlisting}
    which in turn includes -
    \begin{lstlisting}[gobble=4, language=threepl]
       module "xilinx/xc7.3pl"
       source "xilinx/zynq_axi.3pl"
    \end{lstlisting}
    
    Module zynq\_ps() is called as a trailer module (defined in
    xilinx/zynq\_axi.3pl). It instantiates the Zynq CPU and connects up the DDR
    RAM and some general purpose I/O pins. It also derives a 100MHz clock from
    the CPU which is the bus clock and is available as {\bf c100}.
    
    The following are defined as I/O pins with standard "LVCMOS33" -
    
    Test pins tp6\_pin to tp20\_pin("str").
    
    The 50 pins on each of connector 1 and connector 2,
    con1\_pins and con2\_pins ("[50]str").
    
    The I/O pins on each connector are grouped with GND and VCC pins between
    groups. For convenience the pin groups are also defined separately
    as con1\_3\_16\_pins, con1\_19\_38\_pins etc., each being an array of "str".
    




    \subsection{library modules}
    
        None.




    \subsection{library procedures}
    
        None.
        



    \subsection{library functions}
    
        None.


    \subsection{library classes}
    
        None.
